\section{Three-label states: baryon labels in \texorpdfstring{$E_8$}{E8} and in the spinor class}
\label{sec:threelabel}
The adjoint of $SO(32)$ organizes two-label states. This section records
what $E_8$ holds instead, and where three-label states sit on the
$SO(32)$ side. Details and the exact-arithmetic script are in
blackboard~014. The representation content is computed; its reading in
terms of flavor, color and charge is conditional, and statements about
known string models and dualities are quoted from memory.

\paragraph{$E_8$ in $SU(9)$ form.}
The roots of $E_8$ are $e_i-e_j$ and
$\pm(e_i+e_j+e_k-\tfrac13\sum e)$ in $\mathbb R^9$, so that
$248=80+84+\overline{84}$ with $84=\Lambda^3(9)$. Reading the nine
labels as three flavors times three colors,
\begin{equation}
84=(10,1)+(8,8)+(1,10),\qquad 80=(8,1)+(1,8)+(8,8).
\end{equation}
The $(10,1)=S^3(3_f)\otimes\Lambda^3(3_c)$ is the label content of the
baryon decuplet. For $(u,d,s)$ the charge
$Q=\mathrm{diag}(2/3,-1/3,-1/3)\otimes1_3$ is traceless and is a
generator of $E_8$; the color singlets of the 84 then have charges
$2,1,0,-1$ with multiplicities $1,2,3,4$, and those of the 80 have
charges $-1,0,1$ with multiplicities $2,4,2$. The decomposition is a
$\mathbb Z_3$ grading: two elements of the 84 bracket into the
$\overline{84}$, so label number is conserved modulo 9. There is no
$\Lambda^2(9)$, hence no diquark-like state, and no baryon octet, which
would need spin labels and $SU(18)$.

\paragraph{A string model.}
Consider the $\mathbb Z_3$ orbifold of the $E_8\times E_8$ heterotic
string with shift
\begin{equation}
V=\tfrac13\,(1,1,1,1,2,0,0,0\,;\,2,0,0,0,0,0,0,0).
\end{equation}
Its gauge group is $SU(9)\times SO(14)\times U(1)$, its untwisted matter
is $3\,[(84,1)+(1,14)+(1,64)]$ and its twisted matter is
$27\,(\bar9,1)$. The $SU(9)$ cubic anomaly cancels as $3\cdot9-27$. In the label reading the
untwisted matter carries baryon-like labels and the twisted matter
anti-label states, where the $SO(32)$ model of
Section~\ref{sec:gaugefieldscalars} has diquark-like and one-label
states.

\paragraph{The spinor class of $SO(32)$.}
The lattice of $\mathrm{Spin}(32)/\mathbb Z_2$ contains the weights
$(\pm\tfrac12)^{16}$ with an even number of minus signs, that is,
subsets $S$ of the sixteen labels with $|S|$ even. These states are
massive at the string scale. Among them the color singlets have
$0,3,6,9,12,15$ colored labels, with multiplicities
$1,35,175,175,35,1$, and the sixteenth label is present exactly when the
baryon number $B$ is odd. Three consequences follow. Baryon labels exist
in the $SO(32)$ theory, as the $35=S^3(5)$, together with a tower of
multi-baryon label states. The lattice condition reads $3B+L$ even if
the sixteenth label carries lepton number one. And since a subset has
charge $\sum_{a\in S}q_a-\tfrac12\sum_{\rm all}q_a$, with the colored
labels summing to one, the color singlets have integer charge exactly
when the sixteenth label has charge $-1$. This supplies the trace
condition that Section~\ref{sec:gaugefieldscalars} had to assume.

\paragraph{Remarks.}
The decomposition $80+84+\overline{84}$ also organizes the U-duality
algebra $E_{8(8)}\supset SL(9,\mathbb R)$ of M-theory on an eight-torus,
with the 84 holding the three-form components, to which the membrane
couples, and the $\overline{84}$ those of the dual six-form of the
fivebrane; Ref.~\cite{Rivero2009} had matched that 84 with the light
fermion degrees of freedom. This is the same complex algebra in a
different real form and role, and no connection between the two uses is
derived here. An earlier suggestion of this notebook session, that the
$U(9)$ of nine $SO(32)$ labels could enhance to $E_8$ in nine
dimensions, is withdrawn: the duality between the two heterotic frames
does not carry labels to labels, and in the $SO(32)$ frame the
three-label states are those of the spinor class.
